% GENERATED FILE. Edit canonical lessons, then run npm run generate:books.
% canonical-source-hash: fnv1a64:8c2f9bd9c1c368bd
% canonical-lessons: MR-C276-gelo, MR-C276-tyane-kela, MR-C276-hota, MR-R276-what-happened

\chapter{What Happened, Who Did It, What There Was}
\label{ch:mr-what-happened-who-did-it-what-there-was}
\hlchaptermodality{276}

\begin{chapteropening}
\textbf{By the end of this chapter:} \emph{I can say what happened (gelo, ālī), who did something with -ne (tyāne kām kelaṁ), and what there was (hotā).}

Complete MR-R276-what-happened: say what happened, who did it and what there was, in the past.
\end{chapteropening}

\section[\texorpdfstring{\mr{मी} \mr{गेलो} (mī gelo)}{मी गेलो (mī gelo)}]{\mr{मी} \mr{गेलो} (mī gelo) — the past, person by person}

\label{lesson:MR-C276-gelo}



\begin{quote}
Before the new one: say the last two things you named, dance and a stage. (\emph{nṛtya}, \emph{raṅgmañc}.) Then the line you have said since your first meeting: \textbf{\mr{भेटून} \mr{आनंद} \mr{झाला}} (\emph{bheṭūn ānand jhālā}).

Its last word, \textbf{\mr{झाला}}, is a past: \textquotedblleft{}(gladness) came about\textquotedblright{}.
\end{quote}

\subsection*{You'll want to know: the past, person by person}
\begin{quote}
\textbf{\mr{काल} \mr{मी} \mr{पुण्याला} \mr{गेलो}.}
\end{quote}

\begin{quote}
\emph{kāl mī puṇyālā gelo.} — \textquotedblleft{}Yesterday I went to Pune.\textquotedblright{}
\end{quote}

\textbf{\mr{काल}} (\emph{kāl}, \textquotedblleft{}yesterday\textquotedblright{}). \textbf{\mr{गेलो}} is \textbf{\mr{जाणे}} in the past: the stem \textbf{\mr{गे}-}, then \textbf{-\mr{ल}-}, then an ending for person and gender, as in \textbf{\mr{बोलतो} / \mr{बोलते}}:

\begin{itemize}
\raggedright
  \item \textbf{\mr{मी}} \textbf{\mr{गेलो}} (\emph{gelo}) — I went; a woman says \textbf{\mr{मी}} \textbf{\mr{गेले}} (\emph{gele})
  \item \textbf{\mr{तू}} \textbf{\mr{गेलास}} (\emph{gelās}) — you went, to a friend; to a woman, \textbf{\mr{गेलीस}} (\emph{gelīs})
  \item \textbf{\mr{तो}} \textbf{\mr{गेला}} (\emph{gelā}) — he went; \textbf{\mr{ती}} \textbf{\mr{गेली}} (\emph{gelī}) — she went
  \item \textbf{\mr{आम्ही}} (\emph{āmhī}, \textquotedblleft{}we\textquotedblright{}) \textbf{\mr{गेलो}} — we went
  \item \textbf{\mr{तुम्ही}} \textbf{\mr{गेलात}} (\emph{gelāt}) — you went, with respect or to several
  \item \textbf{\mr{ते}} \textbf{\mr{गेले}} (\emph{gele}) — they went; \textbf{\mr{त्या}} (\emph{tyā}, \textquotedblleft{}they, all women\textquotedblright{}) \textbf{\mr{गेल्या}} (\emph{gelyā})
\end{itemize}

\textbf{-\mr{ा}}, \textbf{-\mr{ी}}, \textbf{-\mr{ं}} are the endings of \textbf{\mr{माझा}}, \textbf{\mr{माझी}}, \textbf{\mr{माझं}}. The Marathi past behaves like an adjective that agrees with whoever went.

\subsection*{You'll want to know: the past stems}
Most verbs put \textbf{-\mr{ल}-} straight onto the stem: \textbf{\mr{बसणे}} $\to$ \textbf{\mr{मी} \mr{बसलो}} (\emph{mī baslo}), \textbf{\mr{झोपणे}} $\to$ \textbf{\mr{मी} \mr{झोपलो}} (\emph{mī jhoplo}), \textbf{\mr{पडणे}} $\to$ \textbf{\mr{तो} \mr{पडला}} (\emph{to paḍlā}), \textquotedblleft{}he fell\textquotedblright{}. A few add a vowel first: \textbf{\mr{राहणे}} $\to$ \textbf{\mr{मी} \mr{राहिलो}} (\emph{mī rāhilo}), \textquotedblleft{}I stayed\textquotedblright{}; \textbf{\mr{निघणे}} $\to$ \textbf{\mr{मी} \mr{निघालो}} (\emph{mī nighālo}), \textquotedblleft{}I set out\textquotedblright{}.

Two more change the stem itself:

\begin{itemize}
\raggedright
  \item \textbf{\mr{येणे}} $\to$ \textbf{\mr{मी} \mr{आलो}} (\emph{mī ālo}), \textbf{\mr{ती} \mr{आली}} (\emph{tī ālī}) — came
  \item \textbf{\mr{होणे}} (\emph{hoṇe}, \textquotedblleft{}to become, to happen\textquotedblright{}) $\to$ \textbf{\mr{तो} \mr{झाला}} (\emph{to jhālā}), \textbf{\mr{ते} \mr{झालं}} (\emph{te jhālaṁ}) — became, happened
\end{itemize}

\textquotedblleft{}Did not\textquotedblright{} keeps the past form and puts \textbf{\mr{नाही}} last, as you do with the present: \textbf{\mr{मी} \mr{गेलो} \mr{नाही}} (\emph{mī gelo nāhī}), \textquotedblleft{}I didn't go\textquotedblright{}.

\subsection*{Your turn}
\textbf{\mr{घरी}} (\emph{gharī}, \textquotedblleft{}home, at home\textquotedblright{}).

Say these aloud:

\begin{itemize}
\raggedright
  \item \textquotedblleft{}Yesterday I went to Pune\textquotedblright{} — \emph{kāl mī puṇyālā gelo}
  \item \textquotedblleft{}She came home in the evening\textquotedblright{} — \emph{tī sandhyākāḷī gharī ālī}
  \item \textquotedblleft{}We sat. Then we slept.\textquotedblright{} — \emph{āmhī baslo. mag āmhī jhoplo.}
\end{itemize}

\begin{tcolorbox}[breakable,colback=teal!4,colframe=teal!35!black,title={Before you move on}]
\textquotedblleft{}I went\textquotedblright{}? (\textbf{\mr{गेलो}}, or \textbf{\mr{गेले}}.) \textquotedblleft{}He came\textquotedblright{}? (\textbf{\mr{आला}}.) \textquotedblleft{}It happened\textquotedblright{}? (\textbf{\mr{झालं}}.) Which ending says \textquotedblleft{}she\textquotedblright{}? (\textbf{-\mr{ी}}: \textbf{\mr{गेली}}.) \textquotedblleft{}I didn't go\textquotedblright{}? (\textbf{\mr{मी} \mr{गेलो} \mr{नाही}}.)
\end{tcolorbox}
\section[\texorpdfstring{\mr{त्याने} \mr{काम} \mr{केलं} (tyāne kām kelaṁ)}{त्याने काम केलं (tyāne kām kelaṁ)}]{\mr{त्याने} \mr{काम} \mr{केलं} (tyāne kām kelaṁ) — the past that agrees with the thing}

\label{lesson:MR-C276-tyane-kela}



\begin{quote}
Say \textquotedblleft{}I went\textquotedblright{} and \textquotedblleft{}she came\textquotedblright{}. (\emph{mī gelo}, \emph{tī ālī}.) Then \textquotedblleft{}I see with my eye\textquotedblright{}, with the \textbf{-\mr{ने}} of a tool. (\emph{mī ḍoḷyāne pāhto}.)

That \textbf{-\mr{ने}} was promised a second job. This is it.
\end{quote}

\begin{grammarlens}[title={Grammar Lens: the doer takes -\mr{ने}}]
\begin{quote}
\textbf{\mr{त्याने} \mr{काम} \mr{केलं}.}
\end{quote}

\begin{quote}
\emph{tyāne kām kelaṁ.} — \textquotedblleft{}He did the work.\textquotedblright{}
\end{quote}

A verb with an object, such as eating, reading or doing, changes the past. The doer is marked: \textbf{\mr{तो}} becomes \textbf{\mr{त्याने}} (\emph{tyāne}), \textbf{\mr{ती}} becomes \textbf{\mr{तिने}} (\emph{tine}), and \textbf{\mr{ते}}, \textquotedblleft{}they\textquotedblright{}, becomes \textbf{\mr{त्यांनी}} (\emph{tyānnī}). A name takes \textbf{-\mr{ने}} too: \textbf{\mr{मीराने}} (\emph{mīrāne}).

\textbf{\mr{मी}}, \textbf{\mr{तू}}, \textbf{\mr{आम्ही}} and \textbf{\mr{तुम्ही}} take nothing: \textbf{\mr{मी} \mr{काम} \mr{केलं}} (\emph{mī kām kelaṁ}), \textquotedblleft{}I did the work\textquotedblright{}.
\end{grammarlens}

\subsection*{You'll want to know: the verb follows the thing}
\begin{quote}
\textbf{\mr{मी} \mr{आंबा} \mr{खाल्ला}.}
\end{quote}

\begin{quote}
\emph{mī āmbā khāllā.} — \textquotedblleft{}I ate a mango.\textquotedblright{}
\end{quote}

\textbf{\mr{आंबा}} is masculine, so the verb ends in \textbf{-\mr{ा}}, whoever ate it. Change the thing and the verb changes:

\begin{itemize}
\raggedright
  \item \textbf{\mr{मी} \mr{पोळी}} \textbf{\mr{खाल्ली}} (\emph{khāllī}) — I ate a chapati (feminine)
  \item \textbf{\mr{मी} \mr{पुस्तक}} \textbf{\mr{वाचलं}} (\emph{vāclaṁ}) — I read a book (neuter)
  \item \textbf{\mr{तिने} \mr{चित्रपट}} \textbf{\mr{पाहिला}} (\emph{pāhilā}) — she watched a film
\end{itemize}

After \textbf{\mr{तू}} the verb adds \textbf{-\mr{स}}: \textbf{\mr{तू} \mr{पुस्तक} \mr{वाचलंस}} (\emph{tū pustak vāclaṁs}); after \textbf{\mr{तुम्ही}}, \textbf{-\mr{त}}: \textbf{\mr{तुम्ही} \mr{काय} \mr{केलंत}?} (\emph{tumhī kāy kelant?}), \textquotedblleft{}What did you do?\textquotedblright{}

Pasts to learn: \textbf{\mr{केलं}}, \textbf{\mr{खाल्लं}}, \textbf{\mr{घेतलं}} (\emph{ghetlaṁ}), \textbf{\mr{दिलं}} (\emph{dilaṁ}), \textbf{\mr{सांगितलं}} (\emph{sāṅgitlaṁ}), and from \textbf{\mr{पिणे}}, \textbf{\mr{प्यायलं}} (\emph{pyāylaṁ}). With no object, \textbf{-\mr{ं}}: \textbf{\mr{त्याने} \mr{सांगितलं}} (\emph{tyāne sāṅgitlaṁ}), \textquotedblleft{}he said so\textquotedblright{}. Formal print writes this \textbf{-\mr{ं}} as \textbf{-\mr{े}} (\textbf{\mr{केले}}). A few verbs with an object still follow the doer: \textbf{\mr{मी} \mr{मराठी} \mr{शिकलो}} (\emph{mī marāṭhī śiklo}).

\subsection*{Your turn}
Say these aloud:

\begin{itemize}
\raggedright
  \item \textquotedblleft{}He did the work\textquotedblright{} — \emph{tyāne kām kelaṁ}
  \item \textquotedblleft{}I ate a mango\textquotedblright{} — \emph{mī āmbā khāllā}
  \item \textquotedblleft{}Meera read the news\textquotedblright{} — \emph{mīrāne bātmī vāclī}
\end{itemize}

\begin{tcolorbox}[breakable,colback=teal!4,colframe=teal!35!black,title={Before you move on}]
Who takes \textbf{-\mr{ने}}? (\textbf{\mr{त्याने}}, \textbf{\mr{तिने}}, \textbf{\mr{त्यांनी}}, and a name.) Who takes nothing? (\textbf{\mr{मी}}, \textbf{\mr{तू}}, \textbf{\mr{आम्ही}}, \textbf{\mr{तुम्ही}}.) Why \textbf{\mr{खाल्ली}} in \textbf{\mr{मी} \mr{पोळी} \mr{खाल्ली}}? (\textbf{\mr{पोळी}} is feminine, and the verb follows the thing.) \textquotedblleft{}I read a book\textquotedblright{}? (\textbf{\mr{मी} \mr{पुस्तक} \mr{वाचलं}}.)
\end{tcolorbox}
\section[\texorpdfstring{\mr{होता} (hotā)}{होता (hotā)}]{\mr{होता} (hotā) — was, and there was}

\label{lesson:MR-C276-hota}



\begin{quote}
Say \textquotedblleft{}I am well\textquotedblright{}: \emph{mī barā āhe}. Then \textquotedblleft{}I went\textquotedblright{}: \emph{mī gelo}.

\textbf{\mr{आहे}} has a past of its own, and it takes the endings of \textbf{\mr{गेलो}}.
\end{quote}

\subsection*{You'll want to know: \mr{होता}, 'was'}
\begin{quote}
\textbf{\mr{काल} \mr{मी} \mr{घरी} \mr{होतो}.}
\end{quote}

\begin{quote}
\emph{kāl mī gharī hoto.} — \textquotedblleft{}Yesterday I was at home.\textquotedblright{}
\end{quote}

\begin{itemize}
\raggedright
  \item \textbf{\mr{मी}} \textbf{\mr{होतो}} (\emph{hoto}) — I was; a woman says \textbf{\mr{होते}} (\emph{hote})
  \item \textbf{\mr{तू}} \textbf{\mr{होतास}} (\emph{hotās}) / \textbf{\mr{होतीस}} (\emph{hotīs}) — you were, to a friend
  \item \textbf{\mr{तो}} \textbf{\mr{होता}} (\emph{hotā}), \textbf{\mr{ती}} \textbf{\mr{होती}} (\emph{hotī}), \textbf{\mr{ते}} \textbf{\mr{होतं}} (\emph{hotaṁ}) — he, she, it was
  \item \textbf{\mr{आम्ही}} \textbf{\mr{होतो}}, \textbf{\mr{तुम्ही}} \textbf{\mr{होतात}} (\emph{hotāt}), \textbf{\mr{ते}} \textbf{\mr{होते}} (\emph{hote}) — we, you, they were
\end{itemize}

It is also \textquotedblleft{}there was\textquotedblright{}:

\begin{quote}
\textbf{\mr{गावात} \mr{एक} \mr{नदी} \mr{होती}.}
\end{quote}

\begin{quote}
\emph{gāvāt ek nadī hotī.} — \textquotedblleft{}There was a river in the village.\textquotedblright{}
\end{quote}

\subsection*{You'll want to know: \mr{नव्हता}, 'was not'}
\textquotedblleft{}Was not\textquotedblright{} is one word, with the same endings: \textbf{\mr{नव्हतो}} (\emph{navhto}), \textbf{\mr{नव्हता}} (\emph{navhtā}), \textbf{\mr{नव्हती}} (\emph{navhtī}), \textbf{\mr{नव्हतं}} (\emph{navhtaṁ}).

\begin{quote}
\textbf{\mr{तो} \mr{शाळेत} \mr{नव्हता}.}
\end{quote}

\begin{quote}
\emph{to śāḷet navhtā.} — \textquotedblleft{}He was not at school.\textquotedblright{}
\end{quote}

It is to \textbf{\mr{होता}} what \textbf{\mr{नाही}} is to \textbf{\mr{आहे}}.

\subsection*{Your turn}
Say these aloud:

\begin{itemize}
\raggedright
  \item \textquotedblleft{}Yesterday I was at home\textquotedblright{} — \emph{kāl mī gharī hoto}
  \item \textquotedblleft{}There was a big shop in the village\textquotedblright{} — \emph{gāvāt ek moṭhaṁ dukān hotaṁ}
  \item \textquotedblleft{}He was not at school\textquotedblright{} — \emph{to śāḷet navhtā}
\end{itemize}

\begin{tcolorbox}[breakable,colback=teal!4,colframe=teal!35!black,title={Before you move on}]
\textquotedblleft{}I was\textquotedblright{}? (\textbf{\mr{होतो}}, or \textbf{\mr{होते}}.) \textquotedblleft{}It was\textquotedblright{}? (\textbf{\mr{होतं}}.) \textquotedblleft{}There was a river\textquotedblright{}? (\textbf{\mr{एक} \mr{नदी} \mr{होती}}.) \textquotedblleft{}She was not\textquotedblright{}? (\textbf{\mr{ती} \mr{नव्हती}}.)
\end{tcolorbox}
\section[(review)]{Review — what happened}

\label{lesson:MR-R276-what-happened}



\begin{quote}
Say \textquotedblleft{}Yesterday I went to Pune\textquotedblright{}. (\emph{kāl mī puṇyālā gelo}.) Then \textquotedblleft{}she came\textquotedblright{} and \textquotedblleft{}we sat\textquotedblright{}. (\emph{tī ālī}, \emph{āmhī baslo}.)
\end{quote}

\subsection*{Your turn}
Say these aloud:

\begin{itemize}
\raggedright
  \item I came home in the evening, and then I slept.
  \item Meera ate a mango. Arun read a book.
  \item What did you do yesterday? I was at home.
  \item There was a big shop in the village.
  \item I didn't go to school. He was not at home either.
\end{itemize}

\begin{tcolorbox}[breakable,colback=teal!4,colframe=teal!35!black,title={Before you move on}]
Which past follows the one who went? (\textbf{\mr{गेलो}}, \textbf{\mr{आली}}: no object.) Which follows the thing? (The past with an object: \textbf{\mr{मीराने} \mr{आंबा} \mr{खाल्ला}}.) \textquotedblleft{}There was\textquotedblright{}? (\textbf{\mr{होता}}, \textbf{\mr{होती}}, \textbf{\mr{होतं}}.) \textquotedblleft{}Was not\textquotedblright{}? (\textbf{\mr{नव्हता}}.)
\end{tcolorbox}
